% !TEX root = 02.tex
\documentclass[a4paper,10pt]{article}

\usepackage{pdfsync}
\usepackage[utf8]{inputenc}
\usepackage{microtype}
\usepackage{mathtools}
\usepackage{amsthm,amsfonts,amsmath,amssymb}
\usepackage{mathabx}
% \usepackage{MnSymbol}
\usepackage{hyperref}
% \usepackage{bbold}
\usepackage[textwidth=2cm,textsize=small]{todonotes}
\usepackage{subcaption}
\usepackage{enumitem}
\usepackage{xspace}
% \usepackage{stmaryrd}
\usepackage{verbatim}
\usepackage{cleveref}
\usepackage{csquotes}
\usepackage{cite}

\bibliographystyle{plain}

\hypersetup{
    colorlinks,
    linkcolor={red!50!black},
    citecolor={blue!50!black},
    urlcolor={blue!80!black}
}

\author{}

\newcommand{\ignore}[1]{}

\newcommand{\eg}{\textit{e.g.}\xspace}
\newcommand{\ie}{\textit{i.e.}\xspace}
\newcommand{\ea}{\textit{et al.}\xspace}
\newcommand{\etc}{\textit{etc.}\xspace}
\newcommand{\resp}{resp.\xspace}
\newcommand{\st}{s.t.\xspace}

\newcommand{\sep}{\;|\;}
\renewcommand{\rule}[2]{\dfrac{#1}{#2}}
\newcommand{\set}[1]{\left\{#1\right\}}
\newcommand{\setof}[2]{\set{#1 \;\middle|\; #2}}
\newcommand{\tuple}[1]{\left( #1 \right)}
\newcommand{\pair}[2]{\tuple {#1, #2}}
\newcommand{\closure}[3]{\tuple {#1, #2, #3}}
\newcommand{\sem}[1]{\llbracket{#1}\rrbracket}
\newcommand{\powerset}[1]{\mathcal P(#1)}
\newcommand{\map}[2]{\mathsf{map}\; #1\; #2}
\newcommand{\filter}[2]{\mathsf{filter}\; #1\; #2}
\newcommand{\floor}[1]{\left\lfloor #1 \right\rfloor}
\newcommand{\upwardclosure}[1]{#1\uparrow}
\newcommand{\card}[1]{\left| #1 \right|}
\newcommand{\lang}[1]{L(#1)}
\newcommand{\goesto}[1]{\xrightarrow{#1}}
\newcommand{\e}{\varepsilon}
\newcommand{\supp}[1]{\mathsf{supp}(#1)}
\newcommand{\poly}[2]{#1[#2]}
\newcommand{\rat}[2]{#1(#2)}
\newcommand{\powerseries}[2]{#1[\![#2]\!]}
\newcommand{\dollar}{\$}
\newcommand{\coef}[2]{[#1]#2}

\newcommand{\C}{\mathbb C}
\newcommand{\N}{\mathbb N}
\newcommand{\Z}{\mathbb Z}
\newcommand{\Q}{\mathbb Q}
\newcommand{\K}{\mathbb K}
\newcommand{\R}{\mathbb R}

\newcommand{\NL}{\textsf{NL}}
\newcommand{\NC}[1]{$\textsf{NC}^{#1}$}
\newcommand{\PTIME}{\textsf{PTIME}}
\newcommand{\PSPACE}{\textsf{PSPACE}}

\newtheorem{exercise}{Exercise}

\newenvironment{solution}{\color{teal} \begin{proof}[\protect{\textnormal{\emph{Solution.}}}]}{\end{proof}}

\let\oldqedhere\qedhere

\newif\ifstartedinmathmode
\renewcommand*{\qedhere}{%
  \relax\ifmmode\startedinmathmodetrue\else\startedinmathmodefalse\fi%
  \ifstartedinmathmode\tag*{\protect\oldqedhere}\else\ifinlist{\hfill\oldqedhere}{\oldqedhere}\fi%
}

\crefname{equation}{}{}
\crefname{exercise}{Exercise}{Exercises}


\title{Formal power series and applications - Tutorial 02}
\date{2026/10/13}
\author{Yahia Idriss Benalioua \\ \small University of Warsaw}

\begin{document}

\maketitle

\begin{exercise}[Minimal polynomials]
    \begin{enumerate}
        \item[]
        \item Show that an algebraic formal power series is not
            always uniquely determined by its minimal polynomial.
        \item Suppose the minimal polynomial $P \in
            \poly{\rat{\Q}{x}}{y}$ of an algebraic series $f$ has
            more than one formal power series root. Can there be
            another polynomial $Q \in  \poly{\rat{\Q}{x}}{y}$ that
            uniquely characterizes $f$?
        \item Give examples of polynomials that admit exactly one
            formal power series root, and examples of polynomials
            that admit none.
    \end{enumerate}
\end{exercise}
\begin{solution}
    \begin{enumerate}
        \item[]
        \item Let $f = \sqrt{1 - 4x} = 1 - \sum_{n=1}^{\infty} \frac{2}{n}
            \binom{2n-2}{n-1} x^n$ (see Exercise~\ref{ex:catalan}).
            It is an algebraic series verifying $f^2 - 1 - 4x = 0$,
            but its minimal polynomial $P = y^2 - 1 - 4x$ also admits
            $-f \neq f$ as a root.
        \item That's impossible because, by minimality of $P$, any
            polynomial $Q \in  \poly{\rat{\Q}{x}}{y}$ vanishing on
            $f$ must be a multiple of $P$.
            Consequently, any formal power series root of $P$ must
            also be a root of $Q$.
        \item Any polynomial verifying the conditions of the implicit
            power series theorem must have a unique power series root.
            Take \eg $P = y^2 - xy - y$, we have $P(0,0) = 0$ and
            $\frac{\partial P}{\partial y} = 2y - x - 1$, thus
            $\frac{\partial P}{\partial y} (0,0) = -1 \neq 0$ and $P$
            has a unique power series root.

            Let $P=y^2 - x$.
            If $P$ admitted a power series root $f$, then $f^2 = x$,
            meaning that $1 = \ord{x} = \ord{f^2} = 2 \ord{f}$, which
            is impossible.
    \end{enumerate}
\end{solution}

\begin{exercise}[Composition of formal power series]
    Let $f, g \in \powerseries \Q x$.
    \begin{enumerate}
        \item Under which conditions does the composition $f \circ g$ exist?
        \item Show that rational and algebraic series are closed
            under composition.
    \end{enumerate}
\end{exercise}
\begin{solution}
    \begin{enumerate}
        \item[]
        \item We can show that $f \circ g$ exists if and only if $g(0) = 0$.

            If $g(0) \neq 0$, then the constant term for $f \circ g$
            would be an infinite sum, which is not always defined in general.

            Conversely, if $g(0) = 0$, then $\ord{g^n} \geq n$ for
            all $n \in \N$ and all the coefficients of $f \circ g$
            are the defined by finite sums.

            See~\cite[Section 2.3]{KauersPaule:Tetrahedron:2011} or
            \cite[Section 1.1]{Stanley:EC:CUP:2011} for more formal details.
        \item Let $f = \frac{P}{Q}$ and $g = \frac{R}{S}$ with
            $P,Q,R,S \in \poly{\Q}{x}$ be two rational series.
            We have $f \circ g = \frac{P(R/S)}{Q(R/S)}$, which is
            well-defined since $Q(R/S)(0) = Q(g(0)) = Q(0) \neq 0$ as
            a direct consequence of the conditions for the definition
            of $f$ and $f \circ g$.
            Moreover, it is clearly rational since it can be
            rewritten as the quotient of two polynomials.

            For algebraic series, we can use the characterization
            that a power series $f$ is algebraic if and only if there
            exists a finite dimensional algebra $A$ such that $f \in
            A$, to show the result as follows.

            Let $f, g \in \powerseries{\Q}{x}$ such that $f^d =
            \sum_{i = 0}^{d-1} p_i(x) f^i$ and $ g^r = \sum_{i =
            0}^{r-1} q_i(x) g^i$ for some $p_0, \dots, p_{d-1}, q_0,
            \dots, q_{r-1} \in \rat{\Q}{x}$.
            Let
            \[
                A = \lin[\rat{\Q}{x}]\left( \setof{(f^i \circ g)g^j}{0 \leq
                i \leq d-1 \text{ and } 0 \leq i \leq r-1} \right)
            \]
            Obviously, $f \circ g \in A$, and it remains to show that
            $A$ is closed by products.
            Clearly, it is enough to show that $f^d \circ g \in A$
            and $g^r \in A$.
            For $g^r$, by definition $g^r \in \lin[\rat{\Q}{x}](g^0,
            \dots, g^{r-1}) \subseteq A$.
            As for $f^d \circ g$, we have
            \[
                f^d \circ g
                = \left(\sum_{i = 0}^{d-1} p_i(x) f^i\right) \circ g
                = \sum_{i = 0}^{d-1} p_i(g) f^i \circ g
            \]
            which is an element of $A$ since the $p_i(g) \in
            \lin[\rat{\Q}{x}](g^0,
            \dots, g^{r-1})$.
    \end{enumerate}
\end{solution}

\begin{exercise}[Compositional inverses]
    Let $f \in \powerseries \Q x$ be a power series.
    It admits a \emph{right compositional inverse} if there exists $g
    \in \powerseries \Q x$ such that $f(g(x)) = x$.
    \begin{enumerate}
        \item When does $f$ admit a right compositional inverse?
        \item If $f$ admits a right compositional inverse, does it
            admit a left compositional inverse?
        \item If $f$ admits a right compositional inverse and is
            algebraic, is the inverse also algebraic?
    \end{enumerate}
\end{exercise}
\begin{solution}
    Let us first make some useful observations.
    Let $f = \sum_{n \in \N} f_n x^n$ and let $g = \sum_{n \in \N}
    g_n x^n$ with $g(0) = g_0 = 0$.
    For all $n \in \N$, $\ord{g} \geq n$.
    Thus $\coef{x^n}{f \circ g} = \coef{x^n}{\left(\sum_{k = 0}^{n}
    f_k g^k\right)} = \sum_{k = 0}^{n} f_k \coef{x^n}{g^k}$.
    Moreover, $\coef{x^n}{g^k} = \sum_{i_1 + \dots + i_k = n} g_{i_1}
    \dots g_{i_k}$.
    Thus,
    \begin{equation}\label{eq:coef}
        \coef{x^n}{f \circ g} = \sum_{k = 0}^{n} f_k \sum_{i_1 +
        \dots + i_k = n} g_{i_1} \dots g_{i_k}
    \end{equation}
    \begin{enumerate}
        \item $f$ admits a right compositional inverse if and only if
            $f(0)=0$ and $\frac{\mathrm{d} f}{\mathrm{d} x}(0) \neq 0$.

            Indeed, if there exists $g \in \powerseries{\Q}{x}$ such
            that $f \circ g = x$, then $f \circ g (0) = f(0) = 0$
            and, by \eqref{eq:coef}, $\coef{x}{f \circ g} = f_1 g_1 =
            \frac{\mathrm{d} f}{\mathrm{d} x}(0) \frac{\mathrm{d}
            g}{\mathrm{d} x}(0) = 1$ and thus $\frac{\mathrm{d}
            f}{\mathrm{d} x}(0) \neq 0$.

            Conversely, if $f(0)=0$ and $\frac{\mathrm{d}
            f}{\mathrm{d} x}(0) \neq 0$, then we can define $g$ such
            that $f \circ g = x$ inductively as follows.
            For the composition to be well defined we must have $g(0)=0$.
            Moreover, by identifying the coefficients of $f \circ g$
            and $x$, we have $g_1 = \frac{1}{f_1}$, and $\coef{x^n}{f
            \circ g} = 0$ for all $n \geq 2$.
            Thus, by \eqref{eq:coef}, $g_n = \frac{-1}{f_1} \sum_{k =
            2}^{n} f_k \sum_{i_1 + \dots + i_k = n} g_{i_1} \dots
            g_{i_k}$ is defined as a function of the values of the
            previous $g_i$.
        \item If $f$ admits a right compositional inverse $g$, then,
            by the previous question, $g$ verifies the same
            conditions and admits a right compositional inverse $h$.
            But, by associativity, $h = (f \circ g) \circ h = f \circ
            (g \circ h) = f$, and $g$ is also the left compositional
            inverse of $f$.

            Thus, the formal power series verifying the conditions
            described in the previous question form a group whose
            operation is composition.
            The compositional inverse of $f$ is usually denoted by
            $f^{\left<-1\right>}$.
        \item If $f$ is algebraic then there exists $P \in
            \poly{\Q}{x,y}$ such that $P(x,f)=0$.
            If $f^{\left<-1\right>}$ is the compositional inverse of
            $f$, then $0 = P(x,f) = P(f^{\left<-1\right>},
            f^{\left<-1\right>} \circ f) = P(f^{\left<-1\right>},x)$
            and $f^{\left<-1\right>}$ is thus algebraic.
    \end{enumerate}
\end{solution}

% \begin{exercise}[Lagrange inversion formula]
%     The purpose of this exercise is to develop a purely algebraic
%     proof of the Lagrange inversion formula.
%     \begin{enumerate}
%         \item Let $\phi \in \powerseries \Q w$ be a power series with
%             $\phi(0) \neq 0$ and let $f \in \powerseries \Q x$ be the
%             unique power series such that $f = x \phi(f)$.
%             Show that for all $n \geq 1$,
%             $\coef {x^n} {f(x)} = \frac 1 n \coef {w^{n-1}} {\phi(w)^n}$.
%         \item Let $f$ be unique power series such that $f = x e^f$.
%             Use the Lagrange inversion formula to show that
%             $\coef {x^n} f = \frac {n^{n-1}} {n!}$.
%     \end{enumerate}
% \end{exercise}

\begin{exercise}[Algebraic series and context-free
    languages]\label{ex:catalan}
    Determine the number of words of length $2n$ in the Dyck
    language, \ie the language of well-parenthesized words.

    \emph{Hint:} you may use the identity $\sqrt{1 - 4x} = 1 -
    \sum_{n=1}^{\infty} \frac{2}{n} \binom{2n-2}{n-1} x^n$ to avoid
    tedious calculations (this can be shown using the generalized
    binomial theorem).
\end{exercise}
\begin{solution}
    Let $f = \sum_{n \in \N} f_n x^n$ where $f_n$ is the number of
    words of length $2n$ in the Dyck language.
    Consider the unambiguous context-free grammar defined by its
    rules $S \to \epsilon | (S)S$ generating the Dyck language.

    By translating the rules of the grammar, we obtain $f = 1 + x f^2$.
    Note that, since we're counting pairs of parentheses, the second
    rule yields $x f^2$ instead of $x^2 f^2$.

    Solving the quadratic equation, we obtain either $f = \frac{1 +
    \sqrt{1 - 4x}}{2x}$ or $f = \frac{1 - \sqrt{1 - 4x}}{2x}$.
    The positive root is excluded since $f(0) = 1$.
    Thus
    \begin{align*}
        f &= \frac{1 - \sqrt{1 - 4x}}{2x} \\
        &= \frac{1}{2x} \sum_{n=1}^{\infty} \frac{2}{n} \binom{2n-2}{n-1} x^n \\
        &= \sum_{n=1}^{\infty} \frac{1}{n} \binom{2n-2}{n-1} x^{n-1} \\
        f &= \sum_{n \in \N}\frac{1}{n+1} \binom{2n}{n} x^n
    \end{align*}

    Finally, $f_n = \frac{1}{n+1} \binom{2n}{n}$, which is the $n$-th
    Catalan number.

    See~\cite[Chapter 6]{Sta2023} for a (very long) list of combinatorial
    structures counted by the Catalan numbers.
\end{solution}

\begin{exercise}[Showing transcendence]
    \begin{enumerate}
        \item[]
        \item Show that $e^x$ is not an algebraic power series.
        \item We will now show Eisenstein's theorem, which can be
            stated as follows: if $f = \sum_{n \in \N} f_n x^n \in
            \powerseries{\Q}{x}$ is algebraic, then there exists a
            positive integer $M$ such that $M^n f_n \in \Z$ for
            all $n \in \N$.
            \begin{enumerate}
                \item Show that for all $f \in \powerseries{\Q}{x}$
                    with $f(0) = 0$, if there exists $Q \in
                    \poly{\Z}{x}$ such that, $Q(x,f)=0$, $Q(0,0)=0$
                    and $\frac{\partial Q}{\partial y}(0,0) = c \neq
                    0$ then $f(c^2x) \in \powerseries{\Z}{x}$.
                \item Prove Eisenstein's theorem.
            \end{enumerate}
        \item Deduce an alternative proof of the transcendence of $e^x$.
        \item Is the converse of Eisenstein's theorem true?
    \end{enumerate}
\end{exercise}
\begin{solution}
    \begin{enumerate}
        \item[]
        \item By contradiction, suppose that $e^x$ is algebraic, then
            there exists a minimal $d \in \N$ such that
            \begin{equation}\label{eq:expo}
                \sum_{i = 0}^{d} p_i(x) e^{ix} = 0
            \end{equation}
            for some $p_0, \dots, p_d \in \rat{\Q}{x}$.
            Without loss of generality, we can assume that $p_d = 1$.
            By derivation, we obtain
            \begin{equation}\label{eq:drvexpo}
                \sum_{i = 0}^{d} p_i'(x) e^{ix} + i p_i(x) e^{ix} = 0
            \end{equation}

            Subtracting $d$ times the equation~\eqref{eq:expo}
            from~\eqref{eq:drvexpo} we obtain
            \[
                \sum_{i = 0}^{d-1} (p_i'(x) + (i-d) p_i(x)) e^{ix} = 0
            \]
            contradicting the minimality of $d$.
    \end{enumerate}
    The rest of the questions will be answered in the next tutorial.
\end{solution}

\bibliography{../_bibliography/references}

\end{document}
